\chapter{Fiches de Culture Générale et Entreprise}

Ce chapitre réunit les repères de \textbf{culture générale et entreprise} attendus au QCM : droit des Sociétés Régionales Multiservices, organisation du secteur de l'eau et de l'électricité au Maroc, et actualité récente de la SRM Souss-Massa.

\section{Fiche 1 --- La Loi n° 83-21 en Détail}
\begin{itemize}
    \item \textbf{Intitulé :} loi relative à la gestion des services publics de distribution d'eau potable, d'électricité et d'assainissement liquide.
    \item \textbf{Adoption :} 12 juin 2023, Chambre des Représentants.
    \item \textbf{Promulgation :} dahir n° 1-23-53 du 12 juillet 2023.
    \item \textbf{Publication :} Bulletin Officiel n° 7213 du 17 juillet 2023.
    \item \textbf{Principe central :} confier la gestion de ces trois services, région par région, à une \textbf{société régionale multiservices (SRM)} unique, société anonyme à capitaux publics, se substituant aux régies autonomes municipales et aux concessionnaires privés préexistants.
    \item \textbf{Acteurs remplacés :} régies autonomes (RAMSA à Agadir, RADEEMA à Marrakech, RADEEF à Fès, RADEEJ, etc.) et concessionnaires privés historiques (Lydec à Casablanca, Amendis à Tanger/Tétouan).
    \item \textbf{Calendrier de déploiement :} vague 1 (2024) --- Oriental, Casablanca-Settat, Marrakech-Safi, Souss-Massa ; vagues suivantes (2025-2026) --- Fès-Meknès, Rabat-Salé-Kénitra, Drâa-Tafilalet et autres régions.
\end{itemize}

\section{Fiche 2 --- Le Secteur de l'Eau et de l'Électricité au Maroc}
\begin{itemize}
    \item \textbf{ONEE} (Office National de l'Électricité et de l'Eau potable) : opérateur national historique, producteur d'électricité et d'eau potable en gros, fournisseur des SRM pour une partie de leurs besoins.
    \item \textbf{Stress hydrique national :} le Maroc fait face à une baisse structurelle des ressources en eau conventionnelles (pluviométrie irrégulière, remplissage des barrages en baisse), poussant l'État à investir massivement dans le \textbf{dessalement d'eau de mer} et le transfert d'eau entre bassins.
    \item \textbf{Région Souss-Massa :} l'une des régions les plus exposées au stress hydrique, du fait de la pression agricole (arboriculture, maraîchage d'exportation) et touristique (Agadir) sur la ressource.
    \item \textbf{Dessalement --- exemples nationaux :} Chtouka-Aït Baha (opérationnel, agricole et eau potable), Casablanca (plus grande station de dessalement d'Afrique, en cours de déploiement), Agadir (station historique dédiée à l'irrigation).
\end{itemize}

\section{Fiche 3 --- Chiffres et Dates Clés de la SRM Souss-Massa}
\begin{table}[H]
\centering
\small
\begin{tabular}{|p{5.5cm}|p{8cm}|}
\hline
\rowcolor{gray!20} \textbf{Élément} & \textbf{Valeur} \\ \hline
Loi fondatrice & Loi n° 83-21 \\ \hline
Dahir de promulgation & n° 1-23-53 du 12 juillet 2023 \\ \hline
Publication BO & n° 7213 du 17 juillet 2023 \\ \hline
Entrée en activité & 15 octobre 2024 \\ \hline
Régie précédente & RAMSA (Agadir) \\ \hline
Préfectures couvertes & Agadir Ida-Outanane, Inezgane Aït Melloul \\ \hline
Provinces couvertes & Chtouka-Aït Baha, Taroudant, Tiznit, Tata \\ \hline
Nombre de communes & 175 \\ \hline
Population desservie & environ 3,06 millions d'habitants \\ \hline
Président du Conseil d'Administration & Saaïd Amzazi \\ \hline
Directeur Général & Mohamed Amrzak \\ \hline
Budget d'investissement annoncé & environ 19,68 milliards de dirhams sur 30 ans \\ \hline
Plan quinquennal en cours & 2025--2029 \\ \hline
Postes ouverts au recrutement 2026 & 174 (dont 38 cadres techniques Bac+5) \\ \hline
Code du concours cadres techniques & C43190/26 \\ \hline
\end{tabular}
\caption{Fiche de synthèse à mémoriser avant l'épreuve}
\end{table}

\section{Fiche 4 --- Actualité Récente Utile}
\begin{itemize}
    \item \textbf{27 août 2026 :} annonce du méga-projet de dessalement/transfert d'eau \textbf{Tiznit--Chtouka}, destiné à sécuriser durablement l'approvisionnement en eau potable et agricole de la région Souss-Massa.
    \item \textbf{Barrage Tamri :} ouvrage régional contribuant à la sécurisation de la ressource en eau pour la province d'Agadir Ida-Outanane.
    \item \textbf{Vague 2026 de création de SRM :} poursuite de la réforme avec la création de nouvelles SRM (Fès-Meknès, Casablanca-Settat, Rabat-Salé-Kénitra, Drâa-Tafilalet, etc.), chacune lançant ses propres campagnes de recrutement de cadres techniques.
    \item \textbf{Enjeu de fond :} la réussite de la réforme des SRM est présentée par les pouvoirs publics comme un levier de modernisation de la gouvernance locale des services essentiels, avec une attention particulière portée à la \textbf{digitalisation} (guichet unique, télérelève, SIG) comme facteur de performance.
\end{itemize}

\section{Fiche 5 --- Pièges Fréquents à Éviter}
\begin{itemize}
    \item Ne pas confondre la \textbf{SRM Souss-Massa} avec l'ancienne \textbf{RAMSA} : la RAMSA n'existe plus en tant que structure autonome depuis le 15 octobre 2024, ses missions et son patrimoine ayant été repris par la SRM.
    \item Ne pas confondre le périmètre de la SRM Souss-Massa (2 préfectures + 4 provinces de la région Souss-Massa) avec celui d'autres SRM (chaque région administrative a sa propre SRM).
    \item Le barème exact de l'épreuve écrite (répartition spécialité/culture générale, note éliminatoire) n'étant pas officiellement confirmé, éviter d'affirmer un chiffre précis non vérifié en cas de question ouverte à l'oral ; préférer une formulation prudente (« selon les informations disponibles publiquement »).
\end{itemize}
